% ==============================================================================
% Monografia - Áurea Santos de Oliveira
% Resumo (incluído dentro do ambiente resumo do arquivo principal)
% ==============================================================================
\aluno{Versão preliminar do resumo. Completar com os resultados da avaliação (Capítulo~\ref{sec-avaliacao}) e ajustar para no máximo 500 palavras.}

Sistemas de geração aumentada por recuperação (\textit{Retrieval-Augmented Generation} -- RAG) permitem que modelos de linguagem de grande porte respondam perguntas com base em uma base de conhecimento externa. Na abordagem vetorial tradicional, porém, o texto é recuperado em fragmentos isolados, o que limita respostas que dependem de relações entre partes distantes do corpus -- situação típica de textos legislativos, nos quais artigos remetem uns aos outros e estabelecem definições, exceções e consequências. Este trabalho investiga o GraphRAG, \textit{framework} de código aberto da Microsoft que constrói, com o auxílio de um modelo de linguagem, um grafo de conhecimento a partir do texto, agrupa suas entidades em comunidades hierárquicas e gera relatórios sobre essas comunidades, oferecendo quatro modalidades de consulta: local, global, DRIFT e básica (vetorial). Como prova de conceito, o Código de Defesa do Consumidor (Lei nº 8.078/1990) foi estruturado em um registro por artigo, com sua hierarquia e situação normativa, e indexado com o GraphRAG 3.1.0 utilizando modelos de linguagem em nuvem, uma taxonomia de catorze tipos de entidade definida para o domínio e \textit{prompts} de extração ajustados ao corpus. O índice resultante contém 1.373 entidades, 6.885 relações e 369 comunidades. Foi desenvolvida ainda uma aplicação web, executada em contêineres Docker, que permite explorar o grafo, realizar consultas nas quatro modalidades acompanhando cada etapa da recuperação (contexto recuperado, \textit{prompt} enviado ao modelo, resposta e consumo de \textit{tokens}), destacar no grafo a parte utilizada como contexto e registrar todo o histórico em banco de dados para a avaliação. Durante o desenvolvimento foram identificados e contornados defeitos da biblioteca que afetam diretamente a comparação entre as modalidades, entre eles uma falha que faz a busca básica retornar respostas vazias quando usada programaticamente. \pendente{[Resultados da avaliação comparativa entre as modalidades do GraphRAG e o \textit{baseline} vetorial.]}

\textbf{Palavras-chave}: GraphRAG. Geração aumentada por recuperação. Grafos de conhecimento. Modelos de linguagem. Código de Defesa do Consumidor.
